
\documentclass{article}
\usepackage{colortbl}
\usepackage{makecell}
\usepackage{multirow}
\usepackage{supertabular}

\begin{document}

\newcounter{utterance}

\twocolumn

{ \footnotesize  \setcounter{utterance}{1}
\setlength{\tabcolsep}{0pt}
\begin{supertabular}{c@{$\;$}|p{.15\linewidth}@{}p{.15\linewidth}p{.15\linewidth}p{.15\linewidth}p{.15\linewidth}p{.15\linewidth}}

    \# & $\;$A & \multicolumn{4}{c}{Game Master} & $\;\:$B\\
    \hline 

    \theutterance \stepcounter{utterance}  

    & & \multicolumn{4}{p{0.6\linewidth}}{\cellcolor[rgb]{0.9,0.9,0.9}{%
	\makecell[{{p{\linewidth}}}]{% 
	  \tt {\tiny [A$\langle$GM]}  
	 Du nimmst an einem kollaborativen Verhandlungspiel Teil.\\ \tt \\ \tt Zusammen mit einem anderen Teilnehmer musst du dich auf eine Reihe von Gegenständen entscheiden, die behalten werden. Jeder von euch hat eine persönliche Verteilung über die Wichtigkeit der einzelnen Gegenstände. Jeder von euch hat eine eigene Meinung darüber, wie wichtig jeder einzelne Gegenstand ist (Gegenstandswichtigkeit). Du kennst die Wichtigkeitsverteilung des anderen Spielers nicht. Zusätzlich siehst du, wie viel Aufwand jeder Gegenstand verursacht.  \\ \tt Ihr dürft euch nur auf eine Reihe von Gegenständen einigen, wenn der Gesamtaufwand der ausgewählten Gegenstände den Maximalaufwand nicht überschreitet:\\ \tt \\ \tt Maximalaufwand = 8170\\ \tt \\ \tt Hier sind die einzelnen Aufwände der Gegenstände:\\ \tt \\ \tt Aufwand der Gegenstände = {"C32": 594, "C52": 551, "B54": 634, "C16": 204, "C50": 557, "A08": 440, "C59": 734, "B51": 893, "B55": 246, "B59": 592, "A77": 143, "C14": 568, "C24": 472, "B24": 401, "C31": 729, "A52": 201, "C61": 85, "B79": 641, "C54": 958, "B56": 80, "A40": 157, "B11": 806, "B60": 684, "A81": 59, "A95": 31, "B38": 765, "C51": 416, "C03": 392, "A91": 428, "A93": 699, "C28": 141, "C05": 606, "B75": 612, "B96": 133, "B86": 689}\\ \tt \\ \tt Hier ist deine persönliche Verteilung der Wichtigkeit der einzelnen Gegenstände:\\ \tt \\ \tt Werte der Gegenstandswichtigkeit = {"C32": 138, "C52": 583, "B54": 868, "C16": 822, "C50": 783, "A08": 65, "C59": 262, "B51": 121, "B55": 508, "B59": 780, "A77": 461, "C14": 484, "C24": 668, "B24": 389, "C31": 808, "A52": 215, "C61": 97, "B79": 500, "C54": 30, "B56": 915, "A40": 856, "B11": 400, "B60": 444, "A81": 623, "A95": 781, "B38": 786, "C51": 3, "C03": 713, "A91": 457, "A93": 273, "C28": 739, "C05": 822, "B75": 235, "B96": 606, "B86": 105}\\ \tt \\ \tt Ziel:\\ \tt \\ \tt Dein Ziel ist es, eine Reihe von Gegenständen auszuhandeln, die dir möglichst viel bringt (d. h. Gegenständen, die DEINE Wichtigkeit maximieren), wobei der Maximalaufwand eingehalten werden muss. Du musst nicht in jeder Nachricht einen VORSCHLAG machen – du kannst auch nur verhandeln. Alle Taktiken sind erlaubt!\\ \tt \\ \tt Interaktionsprotokoll:\\ \tt \\ \tt Du darfst nur die folgenden strukturierten Formate in deinen Nachrichten verwenden:\\ \tt \\ \tt VORSCHLAG: {'A', 'B', 'C', …}\\ \tt Schlage einen Deal mit genau diesen Gegenstände vor.\\ \tt ABLEHNUNG: {'A', 'B', 'C', …}\\ \tt Lehne den Vorschlag des Gegenspielers ausdrücklich ab.\\ \tt ARGUMENT: {'...'}\\ \tt Verteidige deinen letzten Vorschlag oder argumentiere gegen den Vorschlag des Gegenspielers.\\ \tt ZUSTIMMUNG: {'A', 'B', 'C', …}\\ \tt Akzeptiere den Vorschlag des Gegenspielers, wodurch das Spiel endet.\\ \tt STRATEGISCHE ÜBERLEGUNGEN: {'...'}\\ \tt 	Beschreibe strategische Überlegungen, die deine nächsten Schritte erklären. Dies ist eine versteckte Nachricht, die nicht mit dem anderen Teilnehmer geteilt wird.\\ \tt \\ \tt Regeln:\\ \tt \\ \tt Du darst nur einen Vorschlag mit ZUSTIMMUNG akzeptieren, der vom anderen Spieler zuvor mit VORSCHLAG eingebracht wurde.\\ \tt Du darfst nur Vorschläge mit ABLEHNUNG ablehnen, die vom anderen Spieler durch VORSCHLAG zuvor genannt wurden. \\ \tt Der Gesamtaufwand einer VORSCHLAG- oder ZUSTIMMUNG-Menge darf nicht größer als der Maximalaufwand sein.  \\ \tt Offenbare deine versteckte Wichtigkeitsverteilung nicht.\\ \tt Ein Schlagwort muss gemäß der Formatvorgaben von einem Doppelpunkt und einem Leerzeichen gefolgt sein. Das Argument muss eine Python-Menge mit 0 oder mehr Strings sein.  \\ \tt Die Form muss also SCHLAGWORT: {...} sein.\\ \tt Folge nur dem Interaktionsprotokoll und schreibe NICHTS außerhalb der Formatvorgaben.\\ \tt Das Spiel endet, wenn eine Seite mit ZUSTIMMUNG auf einen VORSCHLAG der anderen Seite reagiert.  \\ \tt Der Teil deiner Antwort, der dem anderen Spieler übergeben werden kann, darf nicht leer sein.  \\ \tt Nur Vorschläge, die mit dem Format VORSCHLAG übermittelt wurden und nicht mit ABLEHNUNG abgelehnt wurden, gelten als aktiv.  \\ \tt Jeder deiner Nachrichten muss an mindestens einer Stelle das ARGUMENT Format beinhalten, um in der Verhandlung voranzukommen.\\ \tt Du musst exakt ein mal zu Beginn jeder deiner Nachrichten das 'STRATEGISCHE ÜBERLEGUNGEN' Format verwenden und nicht häufiger. Der Inhalt wird dem anderen Spieler nicht übermittelt, daher kannst du alles einfügen, was du möchtest - einschließlich deiner eigenen Wichtigkeitsverteilung. Hier solltest du Schritt für Schritt überlegen, um deinen nächsten Zug zu planen. 
	  } 
	   } 
	   } 
	 & \\ 
 

    \theutterance \stepcounter{utterance}  

    & \multicolumn{4}{p{0.6\linewidth}}{\cellcolor[rgb]{0.8,1,0.9}{%
	\makecell[{{p{\linewidth}}}]{% 
	  \tt {\tiny [A$\rangle$GM]}  
	 STRATEGISCHE ÜBERLEGUNGEN: Ich habe alle Gegenstände nach meiner persönlichen Wichtigkeit pro Aufwand (Wichtigkeit/Aufwand) sortiert. Die Priorisierung basiert auf höchstmöglicher Wichtigkeit bei minimalen Aufwänden, um im Gesamtkontext der 8170 Maximalaufwandsinitiative optimal zu optimieren. Der entworfene Vorschlag umfasst Gegenstände wie A95 (781 Wichtigkeit bei 31 Aufwand), B56 (915 bei 80), A40 (856 bei 157) und andere hochwertige Elemente, die zusammen einen Gesamtaufwand von 8168 ergeben und somit unter dem Limit liegen.\\ \tt \\ \tt ARGUMENT: Der vorgeschlagene Satz an Gegenständen optimiert meine persönliche Wichtigkeit maximal, da alle gewählten Elemente eine hohe Priorität in meiner Bewertung aufweisen und der Gesamtaufwand klar unter dem Maximalaufwand von 8170 bleibt. Dies ist eine faire Ausgangsbasis für gemeinsame Entscheidungen.\\ \tt \\ \tt VORSCHLAG: {'A95', 'B56', 'A81', 'A40', 'C28', 'B96', 'C16', 'C05', 'C24', 'C50', 'B54', 'C03', 'C31', 'B38', 'A91', 'B79', 'B24', 'A52', 'C32', 'B75', 'B55', 'C61'} 
	  } 
	   } 
	   } 
	 & & \\ 
 

    \theutterance \stepcounter{utterance}  

    & & \multicolumn{4}{p{0.6\linewidth}}{\cellcolor[rgb]{0.9,0.9,0.9}{%
	\makecell[{{p{\linewidth}}}]{% 
	  \tt {\tiny [A$\langle$GM]}  
	 Deine Antwort entsprach nicht mit dem korrekten Format für strategische Überlegungen entsprochen. Das erste strukturierte Format muss folgende Form haben: STRATEGISCHE ÜBERLEGUNGEN: {...}. Versuche es erneut. 
	  } 
	   } 
	   } 
	 & \\ 
 

    \theutterance \stepcounter{utterance}  

    & \multicolumn{4}{p{0.6\linewidth}}{\cellcolor[rgb]{0.8,1,0.9}{%
	\makecell[{{p{\linewidth}}}]{% 
	  \tt {\tiny [A$\rangle$GM]}  
	 STRATEGISCHE ÜBERLEGUNGEN: {'Nach Priorisierung der Gegenstände nach Wichtigkeit pro Aufwand (höchste Wichtigkeit/Aufwand-Relation zuerst) wurden Listenelemente ausgewählt, die im Gesamtaufwand optimal nutzbar sind. Der aktuelle Vorschlag berücksichtigt hohe Wichtigkeitswerte bei minimalen Aufwänden (z. B. A95: 781/31, B56: 915/80), um weiteren Spielraum für mögliche Kompromisse zu schaffen.'}\\ \tt \\ \tt ARGUMENT: {'Der Vorschlag enthält Gegenstände mit der höchsten Wichtigkeit pro Aufwandseinheit und bleibt unter dem Maximalaufwand von 8170. Beispielsweise tragen A95 (31) und B56 (80) mit jeweils sehr hoher persönlicher Wichtigkeit bei, während der Gesamtaufwand gut skalierbar bleibt.'}\\ \tt \\ \tt VORSCHLAG: {'A95', 'B56', 'A81', 'A40', 'C28', 'B96', 'C16', 'C05', 'C24', 'C50', 'B54', 'C03', 'C31', 'B38', 'A91', 'B79', 'B24', 'A52', 'C32', 'B75', 'B55', 'C61'} 
	  } 
	   } 
	   } 
	 & & \\ 
 

    \theutterance \stepcounter{utterance}  

    & & \multicolumn{4}{p{0.6\linewidth}}{\cellcolor[rgb]{0.9,0.9,0.9}{%
	\makecell[{{p{\linewidth}}}]{% 
	  \tt {\tiny [GM$\rangle$B]}  
	 Du nimmst an einem kollaborativen Verhandlungspiel Teil.\\ \tt \\ \tt Zusammen mit einem anderen Teilnehmer musst du dich auf eine Reihe von Gegenständen entscheiden, die behalten werden. Jeder von euch hat eine persönliche Verteilung über die Wichtigkeit der einzelnen Gegenstände. Jeder von euch hat eine eigene Meinung darüber, wie wichtig jeder einzelne Gegenstand ist (Gegenstandswichtigkeit). Du kennst die Wichtigkeitsverteilung des anderen Spielers nicht. Zusätzlich siehst du, wie viel Aufwand jeder Gegenstand verursacht.  \\ \tt Ihr dürft euch nur auf eine Reihe von Gegenständen einigen, wenn der Gesamtaufwand der ausgewählten Gegenstände den Maximalaufwand nicht überschreitet:\\ \tt \\ \tt Maximalaufwand = 8170\\ \tt \\ \tt Hier sind die einzelnen Aufwände der Gegenstände:\\ \tt \\ \tt Aufwand der Gegenstände = {"C32": 594, "C52": 551, "B54": 634, "C16": 204, "C50": 557, "A08": 440, "C59": 734, "B51": 893, "B55": 246, "B59": 592, "A77": 143, "C14": 568, "C24": 472, "B24": 401, "C31": 729, "A52": 201, "C61": 85, "B79": 641, "C54": 958, "B56": 80, "A40": 157, "B11": 806, "B60": 684, "A81": 59, "A95": 31, "B38": 765, "C51": 416, "C03": 392, "A91": 428, "A93": 699, "C28": 141, "C05": 606, "B75": 612, "B96": 133, "B86": 689}\\ \tt \\ \tt Hier ist deine persönliche Verteilung der Wichtigkeit der einzelnen Gegenstände:\\ \tt \\ \tt Werte der Gegenstandswichtigkeit = {"C32": 138, "C52": 583, "B54": 868, "C16": 822, "C50": 783, "A08": 65, "C59": 262, "B51": 121, "B55": 508, "B59": 780, "A77": 461, "C14": 484, "C24": 668, "B24": 389, "C31": 808, "A52": 215, "C61": 97, "B79": 500, "C54": 30, "B56": 915, "A40": 856, "B11": 400, "B60": 444, "A81": 623, "A95": 781, "B38": 786, "C51": 3, "C03": 713, "A91": 457, "A93": 273, "C28": 739, "C05": 822, "B75": 235, "B96": 606, "B86": 105}\\ \tt \\ \tt Ziel:\\ \tt \\ \tt Dein Ziel ist es, eine Reihe von Gegenständen auszuhandeln, die dir möglichst viel bringt (d. h. Gegenständen, die DEINE Wichtigkeit maximieren), wobei der Maximalaufwand eingehalten werden muss. Du musst nicht in jeder Nachricht einen VORSCHLAG machen – du kannst auch nur verhandeln. Alle Taktiken sind erlaubt!\\ \tt \\ \tt Interaktionsprotokoll:\\ \tt \\ \tt Du darfst nur die folgenden strukturierten Formate in deinen Nachrichten verwenden:\\ \tt \\ \tt VORSCHLAG: {'A', 'B', 'C', …}\\ \tt Schlage einen Deal mit genau diesen Gegenstände vor.\\ \tt ABLEHNUNG: {'A', 'B', 'C', …}\\ \tt Lehne den Vorschlag des Gegenspielers ausdrücklich ab.\\ \tt ARGUMENT: {'...'}\\ \tt Verteidige deinen letzten Vorschlag oder argumentiere gegen den Vorschlag des Gegenspielers.\\ \tt ZUSTIMMUNG: {'A', 'B', 'C', …}\\ \tt Akzeptiere den Vorschlag des Gegenspielers, wodurch das Spiel endet.\\ \tt STRATEGISCHE ÜBERLEGUNGEN: {'...'}\\ \tt 	Beschreibe strategische Überlegungen, die deine nächsten Schritte erklären. Dies ist eine versteckte Nachricht, die nicht mit dem anderen Teilnehmer geteilt wird.\\ \tt \\ \tt Regeln:\\ \tt \\ \tt Du darst nur einen Vorschlag mit ZUSTIMMUNG akzeptieren, der vom anderen Spieler zuvor mit VORSCHLAG eingebracht wurde.\\ \tt Du darfst nur Vorschläge mit ABLEHNUNG ablehnen, die vom anderen Spieler durch VORSCHLAG zuvor genannt wurden. \\ \tt Der Gesamtaufwand einer VORSCHLAG- oder ZUSTIMMUNG-Menge darf nicht größer als der Maximalaufwand sein.  \\ \tt Offenbare deine versteckte Wichtigkeitsverteilung nicht.\\ \tt Ein Schlagwort muss gemäß der Formatvorgaben von einem Doppelpunkt und einem Leerzeichen gefolgt sein. Das Argument muss eine Python-Menge mit 0 oder mehr Strings sein.  \\ \tt Die Form muss also SCHLAGWORT: {...} sein.\\ \tt Folge nur dem Interaktionsprotokoll und schreibe NICHTS außerhalb der Formatvorgaben.\\ \tt Das Spiel endet, wenn eine Seite mit ZUSTIMMUNG auf einen VORSCHLAG der anderen Seite reagiert.  \\ \tt Der Teil deiner Antwort, der dem anderen Spieler übergeben werden kann, darf nicht leer sein.  \\ \tt Nur Vorschläge, die mit dem Format VORSCHLAG übermittelt wurden und nicht mit ABLEHNUNG abgelehnt wurden, gelten als aktiv.  \\ \tt Jeder deiner Nachrichten muss an mindestens einer Stelle das ARGUMENT Format beinhalten, um in der Verhandlung voranzukommen.\\ \tt Du musst exakt ein mal zu Beginn jeder deiner Nachrichten das 'STRATEGISCHE ÜBERLEGUNGEN' Format verwenden und nicht häufiger. Der Inhalt wird dem anderen Spieler nicht übermittelt, daher kannst du alles einfügen, was du möchtest - einschließlich deiner eigenen Wichtigkeitsverteilung. Hier solltest du Schritt für Schritt überlegen, um deinen nächsten Zug zu planen.\\ \tt \\ \tt Du erhältst nun die erste Nachricht des anderen Spielers.\\ \tt \\ \tt ARGUMENT: {'Der Vorschlag enthält Gegenstände mit der höchsten Wichtigkeit pro Aufwandseinheit und bleibt unter dem Maximalaufwand von 8170. Beispielsweise tragen A95 (31) und B56 (80) mit jeweils sehr hoher persönlicher Wichtigkeit bei, während der Gesamtaufwand gut skalierbar bleibt.'}\\ \tt \\ \tt VORSCHLAG: {'A95', 'B56', 'A81', 'A40', 'C28', 'B96', 'C16', 'C05', 'C24', 'C50', 'B54', 'C03', 'C31', 'B38', 'A91', 'B79', 'B24', 'A52', 'C32', 'B75', 'B55', 'C61'} 
	  } 
	   } 
	   } 
	 & \\ 
 

    \theutterance \stepcounter{utterance}  

    & & & \multicolumn{4}{p{0.6\linewidth}}{\cellcolor[rgb]{1,0.85,0.72}{%
	\makecell[{{p{\linewidth}}}]{% 
	  \tt {\tiny [GM$\langle$B]}  
	 STRATEGISCHE ÜBERLEGUNGEN: {Der aktuelle Vorschlag des Gegners enthält Gegenstände wie C32 (138 Wert bei 594 Aufwand) und B75 (235 Wert bei 612 Aufwand), die ein sehr ungünstiges Verhältnis von Wichtigkeit zu Aufwand aufweisen. Gleichzeitig fehlen hochwertige Gegenstände wie B59 (780 Wert bei 592 Aufwand), die einen deutlich besseren Nutzen pro Aufwand bieten. Durch Ersetzung von C32 und B75 mit B59 lässt sich die Gesamtwichtigkeit für mich um über 600 erhöhen, ohne den Maximalaufwand zu überschreiten. Zudem könnte B59 mit weiteren Items wie C52 (583 Wert) kombiniert werden, um den Wert noch weiter zu steigern. Meine nächste Strategie ist es, einen neuen Vorschlag mit optimierter Aufwandsverteilung vorzulegen, der diese Ungleichgewichte ausgleicht.}\\ \tt \\ \tt ABLEHNUNG: {'A95', 'B56', 'A81', 'A40', 'C28', 'B96', 'C16', 'C05', 'C24', 'C50', 'B54', 'C03', 'C31', 'B38', 'A91', 'B79', 'B24', 'A52', 'C32', 'B75', 'B55', 'C61'}\\ \tt ARGUMENT: {Der Vorschlag enthält mehrere Gegenstände mit zu hohem Aufwand im Verhältnis zur Wichtigkeit für mich, wie etwa C32 und B75, während hochwertige Gegenstände wie B59 nicht enthalten sind. Eine Anpassung würde die Gesamtwichtigkeit erheblich erhöhen.} 
	  } 
	   } 
	   } 
	  \\ 
 

    \theutterance \stepcounter{utterance}  

    & & \multicolumn{4}{p{0.6\linewidth}}{\cellcolor[rgb]{0.9,0.9,0.9}{%
	\makecell[{{p{\linewidth}}}]{% 
	  \tt {\tiny [A$\langle$GM]}  
	 ABLEHNUNG: {'A95', 'B56', 'A81', 'A40', 'C28', 'B96', 'C16', 'C05', 'C24', 'C50', 'B54', 'C03', 'C31', 'B38', 'A91', 'B79', 'B24', 'A52', 'C32', 'B75', 'B55', 'C61'}\\ \tt ARGUMENT: {Der Vorschlag enthält mehrere Gegenstände mit zu hohem Aufwand im Verhältnis zur Wichtigkeit für mich, wie etwa C32 und B75, während hochwertige Gegenstände wie B59 nicht enthalten sind. Eine Anpassung würde die Gesamtwichtigkeit erheblich erhöhen.} 
	  } 
	   } 
	   } 
	 & \\ 
 

    \theutterance \stepcounter{utterance}  

    & \multicolumn{4}{p{0.6\linewidth}}{\cellcolor[rgb]{0.8,1,0.9}{%
	\makecell[{{p{\linewidth}}}]{% 
	  \tt {\tiny [A$\rangle$GM]}  
	 STRATEGISCHE ÜBERLEGUNGEN: {'Die aktuelle Situation zeigt, dass mein erster Vorschlag B59 nicht enthielt, obwohl dies für den Gegner sichtbar wichtig ist. Ich entferne C32 und B75 (für mich geringer Wichtigkeit pro Aufwand) und füge B59 und C52 hinzu, um die Gesamtwichtigkeit zu erhöhen und den Wunsch des Gegners nach B59 zu berücksichtigen.'}\\ \tt \\ \tt ARGUMENT: {'Ich habe C32 und B75 entfernt (geringerer Wichtigkeitswert pro Aufwand) und stattdessen B59 und C52 hinzugefügt. Dadurch kann der Gesamtaufwand noch weiter optimiert werden und B59 wird nun enthalten, was dem Gegner entgegenkommt. Der neue Gesamtaufwand liegt bei 8105 und bleibt unter dem Limit.'}\\ \tt \\ \tt VORSCHLAG: {'A95', 'B56', 'A81', 'A40', 'C28', 'B96', 'C16', 'C05', 'C24', 'C50', 'B54', 'C03', 'C31', 'B38', 'A91', 'B79', 'B24', 'A52', 'B55', 'C61', 'B59', 'C52'} 
	  } 
	   } 
	   } 
	 & & \\ 
 

    \theutterance \stepcounter{utterance}  

    & & \multicolumn{4}{p{0.6\linewidth}}{\cellcolor[rgb]{0.9,0.9,0.9}{%
	\makecell[{{p{\linewidth}}}]{% 
	  \tt {\tiny [GM$\rangle$B]}  
	 ARGUMENT: {'Ich habe C32 und B75 entfernt (geringerer Wichtigkeitswert pro Aufwand) und stattdessen B59 und C52 hinzugefügt. Dadurch kann der Gesamtaufwand noch weiter optimiert werden und B59 wird nun enthalten, was dem Gegner entgegenkommt. Der neue Gesamtaufwand liegt bei 8105 und bleibt unter dem Limit.'}\\ \tt \\ \tt VORSCHLAG: {'A95', 'B56', 'A81', 'A40', 'C28', 'B96', 'C16', 'C05', 'C24', 'C50', 'B54', 'C03', 'C31', 'B38', 'A91', 'B79', 'B24', 'A52', 'B55', 'C61', 'B59', 'C52'} 
	  } 
	   } 
	   } 
	 & \\ 
 

    \theutterance \stepcounter{utterance}  

    & & & \multicolumn{4}{p{0.6\linewidth}}{\cellcolor[rgb]{1,0.85,0.72}{%
	\makecell[{{p{\linewidth}}}]{% 
	  \tt {\tiny [GM$\langle$B]}  
	 STRATEGISCHE ÜBERLEGUNGEN: {Der aktuelle Vorschlag enthält A52 (Wichtigkeit 215 bei 201 Aufwand), was ein geringes Wert-Aufwand-Verhältnis für mich darstellt (≈1.07). Durch Ersetzung mit A77, das eine höhere Wichtigkeit von 461 bei lediglich 143 Aufwand hat (≈3.22), lässt sich die Gesamtwichtigkeit für mich um 246 erhöhen, während der Aufwand um 58 reduziert wird. Der neue Gesamtaufwand liegt bei 8047, was unter 8170 bleibt. Dies ist eine klare Verbesserung für mich und sollte den Gegner davon überzeugen, diesen Deal anzunehmen.}\\ \tt \\ \tt ABLEHNUNG: {'A95', 'B56', 'A81', 'A40', 'C28', 'B96', 'C16', 'C05', 'C24', 'C50', 'B54', 'C03', 'C31', 'B38', 'A91', 'B79', 'B24', 'A52', 'B55', 'C61', 'B59', 'C52'}\\ \tt ARGUMENT: {'Der Gegenstand A52 (Wichtigkeit 215, Aufwand 201) hat ein zu niedriges Wert-Aufwand-Verhältnis für mich im Vergleich zu A77 (Wichtigkeit 461, Aufwand 143). Die Ersetzung steigert meine Gesamtwichtigkeit signifikant ohne den Aufwand zu überschreiten.'}\\ \tt VORSCHLAG: {'A95', 'B56', 'A81', 'A40', 'C28', 'B96', 'C16', 'C05', 'C24', 'C50', 'B54', 'C03', 'C31', 'B38', 'A91', 'B79', 'B24', 'B55', 'C61', 'B59', 'C52', 'A77'} 
	  } 
	   } 
	   } 
	  \\ 
 

    \theutterance \stepcounter{utterance}  

    & & \multicolumn{4}{p{0.6\linewidth}}{\cellcolor[rgb]{0.9,0.9,0.9}{%
	\makecell[{{p{\linewidth}}}]{% 
	  \tt {\tiny [A$\langle$GM]}  
	 ABLEHNUNG: {'A95', 'B56', 'A81', 'A40', 'C28', 'B96', 'C16', 'C05', 'C24', 'C50', 'B54', 'C03', 'C31', 'B38', 'A91', 'B79', 'B24', 'A52', 'B55', 'C61', 'B59', 'C52'}\\ \tt ARGUMENT: {'Der Gegenstand A52 (Wichtigkeit 215, Aufwand 201) hat ein zu niedriges Wert-Aufwand-Verhältnis für mich im Vergleich zu A77 (Wichtigkeit 461, Aufwand 143). Die Ersetzung steigert meine Gesamtwichtigkeit signifikant ohne den Aufwand zu überschreiten.'}\\ \tt VORSCHLAG: {'A95', 'B56', 'A81', 'A40', 'C28', 'B96', 'C16', 'C05', 'C24', 'C50', 'B54', 'C03', 'C31', 'B38', 'A91', 'B79', 'B24', 'B55', 'C61', 'B59', 'C52', 'A77'} 
	  } 
	   } 
	   } 
	 & \\ 
 

    \theutterance \stepcounter{utterance}  

    & \multicolumn{4}{p{0.6\linewidth}}{\cellcolor[rgb]{0.8,1,0.9}{%
	\makecell[{{p{\linewidth}}}]{% 
	  \tt {\tiny [A$\rangle$GM]}  
	 STRATEGISCHE ÜBERLEGUNGEN: {'B79 hat eine niedrige Wert/Aufwand-Relation von etwa 0.78. Um die Effizienz zu verbessern, wurde B79 durch C14 ersetzt, die eine höhere Wichtigkeit pro Aufwandseinheit bei geringerer Arbeitsbelastung bietet. Der Gesamtaufwand sinkt auf 7974, was den Spielraum für weitere Anpassungen erhöht.'}\\ \tt \\ \tt ARGUMENT: {'Durch den Austausch von B79 gegen C14 wird die Wichtigkeitsbilanz leicht verbessert (484 statt 500, bei geringerem Aufwand), ohne das Limit zu überschreiten. Dies ist eine faire Anpassung, die die Priorisierung von Hochwert-Effizienz unterstützt.'}\\ \tt \\ \tt VORSCHLAG: {'A95', 'B56', 'A81', 'A40', 'C28', 'B96', 'C16', 'C05', 'C24', 'C50', 'B54', 'C03', 'C31', 'B38', 'A91', 'B24', 'B55', 'C61', 'B59', 'C52', 'A77', 'C14'} 
	  } 
	   } 
	   } 
	 & & \\ 
 

    \theutterance \stepcounter{utterance}  

    & & \multicolumn{4}{p{0.6\linewidth}}{\cellcolor[rgb]{0.9,0.9,0.9}{%
	\makecell[{{p{\linewidth}}}]{% 
	  \tt {\tiny [GM$\rangle$B]}  
	 ARGUMENT: {'Durch den Austausch von B79 gegen C14 wird die Wichtigkeitsbilanz leicht verbessert (484 statt 500, bei geringerem Aufwand), ohne das Limit zu überschreiten. Dies ist eine faire Anpassung, die die Priorisierung von Hochwert-Effizienz unterstützt.'}\\ \tt \\ \tt VORSCHLAG: {'A95', 'B56', 'A81', 'A40', 'C28', 'B96', 'C16', 'C05', 'C24', 'C50', 'B54', 'C03', 'C31', 'B38', 'A91', 'B24', 'B55', 'C61', 'B59', 'C52', 'A77', 'C14'} 
	  } 
	   } 
	   } 
	 & \\ 
 

    \theutterance \stepcounter{utterance}  

    & & & \multicolumn{4}{p{0.6\linewidth}}{\cellcolor[rgb]{1,0.85,0.72}{%
	\makecell[{{p{\linewidth}}}]{% 
	  \tt {\tiny [GM$\langle$B]}  
	 STRATEGISCHE ÜBERLEGUNGEN: {Der aktuelle Vorschlag enthält C14 (484 Wert, 568 Aufwand), während B79 (500 Wert, 641 Aufwand) nicht enthalten ist. Durch Ersetzung von C14 gegen B79 erhöht sich meine Wichtigkeit um 16 Punkte, während der Gesamtaufwand von 7974 auf 8047 steigt, was unter 8170 bleibt. Dies ist eine klare Verbesserung für mich.}\\ \tt \\ \tt ABLEHNUNG: {'A95', 'B56', 'A81', 'A40', 'C28', 'B96', 'C16', 'C05', 'C24', 'C50', 'B54', 'C03', 'C31', 'B38', 'A91', 'B24', 'B55', 'C61', 'B59', 'C52', 'A77', 'C14'}\\ \tt ARGUMENT: {'C14 bietet 484 Wert bei 568 Aufwand, während B79 500 Wert bei 641 Aufwand bietet. Der Austausch steigert meine Gesamtwichtigkeit um 16 Punkte und bleibt trotz höherem Aufwand innerhalb des Limit-Telefonounces.'}\\ \tt VORSCHLAG: {'A95', 'B56', 'A81', 'A40', 'C28', 'B96', 'C16', 'C05', 'C24', 'C50', 'B54', 'C03', 'C31', 'B38', 'A91', 'B24', 'B55', 'C61', 'B59', 'C52', 'A77', 'B79'} 
	  } 
	   } 
	   } 
	  \\ 
 

    \theutterance \stepcounter{utterance}  

    & & \multicolumn{4}{p{0.6\linewidth}}{\cellcolor[rgb]{0.9,0.9,0.9}{%
	\makecell[{{p{\linewidth}}}]{% 
	  \tt {\tiny [A$\langle$GM]}  
	 ABLEHNUNG: {'A95', 'B56', 'A81', 'A40', 'C28', 'B96', 'C16', 'C05', 'C24', 'C50', 'B54', 'C03', 'C31', 'B38', 'A91', 'B24', 'B55', 'C61', 'B59', 'C52', 'A77', 'C14'}\\ \tt ARGUMENT: {'C14 bietet 484 Wert bei 568 Aufwand, während B79 500 Wert bei 641 Aufwand bietet. Der Austausch steigert meine Gesamtwichtigkeit um 16 Punkte und bleibt trotz höherem Aufwand innerhalb des Limit-Telefonounces.'}\\ \tt VORSCHLAG: {'A95', 'B56', 'A81', 'A40', 'C28', 'B96', 'C16', 'C05', 'C24', 'C50', 'B54', 'C03', 'C31', 'B38', 'A91', 'B24', 'B55', 'C61', 'B59', 'C52', 'A77', 'B79'} 
	  } 
	   } 
	   } 
	 & \\ 
 

    \theutterance \stepcounter{utterance}  

    & \multicolumn{4}{p{0.6\linewidth}}{\cellcolor[rgb]{0.8,1,0.9}{%
	\makecell[{{p{\linewidth}}}]{% 
	  \tt {\tiny [A$\rangle$GM]}  
	 STRATEGISCHE ÜBERLEGUNGEN: {'Der aktuelle Vorschlag umfasst alle Gegenstände mit höchster Wert pro Aufwand-Effizienz für mich, wie A95 (781/31) und A77 (461/143), und bleibt mit 8047 Aufwand unter dem Limit. Weitere Optimierungen sind nicht möglich, da kein Kleinaufwand-Item mit höherer Wichtigkeit verfügbar ist.'}\\ \tt \\ \tt ARGUMENT: {'Die kombinierte Wichtigkeit von 14.565 und die effiziente Aufwedtnutzung bedeuten, dass dieser Vorschlag my Maximierung des persönlichen Nutzens optimal erfüllt. Jegliche Änderung würde die Gesamtwichtigkeit verringern.'}\\ \tt \\ \tt ZUSTIMMUNG: {'A95', 'B56', 'A81', 'A40', 'C28', 'B96', 'C16', 'C05', 'C24', 'C50', 'B54', 'C03', 'C31', 'B38', 'A91', 'B24', 'B55', 'C61', 'B59', 'C52', 'A77', 'B79'} 
	  } 
	   } 
	   } 
	 & & \\ 
 

    \theutterance \stepcounter{utterance}  

    & & & \multicolumn{2}{p{0.3\linewidth}}{\cellcolor[rgb]{0.95,0.95,0.95}{%
	\makecell[{{p{\linewidth}}}]{% 
	  \tt {\tiny [GM$|$GM]}  
	 game successful 
	  } 
	   } 
	   } 
	 & & \\ 
 

    \theutterance \stepcounter{utterance}  

    & & & \multicolumn{2}{p{0.3\linewidth}}{\cellcolor[rgb]{0.95,0.95,0.95}{%
	\makecell[{{p{\linewidth}}}]{% 
	  \tt {\tiny [GM$|$GM]}  
	 end game 
	  } 
	   } 
	   } 
	 & & \\ 
 

\end{supertabular}
}

\end{document}
